\SourceLeaf{176}It is easy to prove that if Euclid's geometry is not the true one,
there are no similar figures at all. The angles in an equilateral triangle
then also vary with the length of the side, in which I find nothing
absurd at all. The angle is then a function of the side, and the side a function
of the angle, naturally a function in which there also occurs a
constant line. It seems somewhat paradoxical that a constant line
  should be possible, as it were, a priori. Yet I find nothing
contradictory in this. It would even be desirable for Euclid's geometry not
  to be true, since we would then have a universal measure a priori. For example,
one could take as the unit of space the side of the equilateral triangle
  whose angle is $=59^{0} 59' 59''{,}99999$. .\ .\ .\ .\ .\ .\ .

\begin{center}\rule{3cm}{0.4pt}\end{center}
